%%%%%%%%%%%%%%%%%%%%%%%%%%%%%%%%%%%%%%%%%
% رزومۀ صنعتی (نسخۀ فارسی)
% معادل فارسیِ ../Industry/Kahbod-Aeini_Resume.tex
% کامپایل: xelatex Kahbod-Aeini_Resume_FA.tex
%%%%%%%%%%%%%%%%%%%%%%%%%%%%%%%%%%%%%%%%%

\documentclass[11pt]{article}
\usepackage{fa-cv}

\begin{document}

\faHeader{کهبد آیینی}{%
واترلو، انتاریو، کانادا\\[0.3em]
\lr{\href{mailto:kaeini@uwaterloo.ca}{kaeini@uwaterloo.ca}} |
\lr{\href{https://www.linkedin.com/in/kahbod-aeini-833b71190/}{LinkedIn}} |
\lr{\href{https://github.com/kahbodaeini}{GitHub}} |
\lr{\href{https://kahbod.org}{kahbod.org}} |
\lr{\href{https://scholar.google.ca/citations?user=3Fd5R9YAAAAJ&hl=en}{Google Scholar}}%
}

\begin{faSection}{خلاصه}
مهندس نرم‌افزار و پژوهشگر سیستم‌های توزیع‌شده. طراحی و تحویل واسط‌های \lr{GraphQL} و \lr{REST}، سرویس‌های متکی بر \lr{SQL} و قابلیت‌های امنیتی برای یک سکوی عملیاتی در حوزۀ خرده‌فروشی انرژی، هم‌زمان با گذراندن دکتری در زمینۀ همروندی و حافظۀ پایدار در دانشگاه واترلو. تسلط بر پشتۀ سمت سرور --- سرویس‌های \lr{C\#/.NET} و \lr{Python}، \lr{SQL} در حجم بالا، پایگاه‌های دادۀ رابطه‌ای و سندی، \lr{Docker} و \lr{Linux} --- به‌همراه تخصص عمیق در ساختمان‌های دادۀ همروند و بدون قفل روی سخت‌افزار چندهسته‌ای واقعی. برندۀ جایزۀ بهترین مقالۀ \lr{PODC 2026} برای کار روی همگام‌سازی بی‌انتظار. هم‌زمان دکتری تمام‌وقت و نقش مهندسی عملیاتی را پیش می‌برد و در هر دو محیط در استخدام، آموزش نیروی تازه و هدایت تیم نقش داشته است.
\end{faSection}

\begin{faSection}{مهارت‌های فنی}

\textbf{زبان‌ها}: \lr{C\#}، \lr{SQL}، \lr{GraphQL}، \lr{Python}، \lr{C/C++}، \lr{Java}، \lr{R}، \lr{Haskell}، \lr{Prolog}

\vspace{0.3em}
\textbf{واسط‌های برنامه‌نویسی و سمت سرور}: طراحی شِما و \lr{resolver} در \lr{GraphQL}، طراحی \lr{REST API}، \lr{.NET / ASP.NET Core}، \lr{Entity Framework}، \lr{Django}، یکپارچه‌سازی \lr{EDI}، \lr{Docker}، \lr{Git}، \lr{Linux / NixOS}

\vspace{0.3em}
\textbf{امنیت}: احراز هویت و مجوزدهی، کنترل دسترسی مبتنی بر نقش و مجوز، رمزنگاری و تبادل امن داده، طراحی امنِ واسط‌های برنامه‌نویسی

\vspace{0.3em}
\textbf{پایگاه‌های داده و داده}: \lr{Microsoft SQL Server}، \lr{PostgreSQL}، \lr{MySQL}، \lr{MongoDB}، پرس‌وجوهای پیچیدۀ \lr{SQL} و بهینه‌سازی آن‌ها، طراحی شِما، طراحی خط لولۀ \lr{ETL}، \lr{Pandas}، \lr{NumPy}

\vspace{0.3em}
\textbf{همروندی و سیستم‌ها}: برنامه‌نویسی چندهسته‌ای و \lr{NUMA}، حافظۀ تراکنشی \lr{Intel RTM}، \lr{OpenMP}، اولیه‌های اتمی (\lr{CAS})، ساختمان‌های دادۀ بدون قفل و بی‌انتظار، بازپس‌گیری حافظۀ مبتنی بر دوره (\lr{EBR})، پروفایل‌گیری با \lr{perf}

\vspace{0.3em}
\textbf{یادگیری ماشین}: \lr{PyTorch}، \lr{TensorFlow}، \lr{scikit-learn}، \lr{OpenCV}، مهندسی ویژگی، ارزیابی مدل

\end{faSection}

\begin{faSection}{سوابق کاری}

\faEntry{مهندس نرم‌افزار --- \lr{ArcTrade} (پاره‌وقت، هم‌زمان با دکتری)}{دسامبر ۲۰۲۵ تاکنون}
\begin{faItems}
  \item طراحی و ساخت واسط‌های \lr{GraphQL} و \lr{REST} برای سامانۀ عامل خرده‌فروشی انرژیِ \lr{ArcTrade}، که شرکت‌های خرده‌فروش برق در بازارهای آمریکای شمالی از آن استفاده می‌کنند.
  \item مشارکت اصلی در ماژول امنیت سکو: احراز هویت، مجوزدهی و کنترل دسترسی مبتنی بر مجوز، که تعیین می‌کند کدام کاربران و شرکا اجازۀ خواندن یا اقدام روی داده‌های حساس بازار و مالی را دارند.
  \item نوشتن و بهینه‌سازی پرس‌وجوهای گستردۀ \lr{SQL} روی مدل دادۀ رابطه‌ای سکو، به‌عنوان پشتوانۀ \lr{resolver}های واسط برنامه‌نویسی و مسیرهای گزارش‌گیری.
  \item پیاده‌سازی تبادل الکترونیکی دادۀ امن و خودکار (\lr{EDI}) با شرکای تجاری بیرونی، برای مبادلۀ رمزنگاری‌شده و تأییدپذیر تراکنش‌های کسب‌وکار.
  \item مشارکت در بخش‌های پشتیبانِ معاملات انرژی و مدیریت ریسک، قیمت‌گذاری، صورت‌حساب و تسویه، جایی که درستی محاسبات مالی جای خطا ندارد.
  \item همکاری در طراحی فرایند استخدام و آموزش نیروهای تازه‌وارد تیم.
\end{faItems}

\faEntry{مهندس داده --- \lr{AI Lida}}{فوریه تا سپتامبر ۲۰۲۲}
\begin{faItems}
  \item توسعه، نگهداری و بهینه‌سازی خطوط لولۀ \lr{ETL} با \lr{Python} و \lr{PostgreSQL} برای بارهای کاری تحلیلی.
  \item همکاری با واحد امنیت برای گردآوری نیازمندی‌ها و ارائۀ راهکارهای دادهایِ هم‌راستا با اهداف کسب‌وکار.
  \item مصاحبه، استخدام، آموزش و سرپرستی تیمی پنج‌نفره از کارآموزان مهندسی داده.
\end{faItems}

\end{faSection}

\begin{faSection}{سوابق پژوهشی}

\faEntry{دستیار پژوهشی تحصیلات تکمیلی --- آزمایشگاه الگوریتم‌ها و سیستم‌های توزیع‌شده، دانشگاه واترلو}{ژانویۀ ۲۰۲۶ تاکنون}
استاد راهنما: پروفسور \lr{Wojciech Golab}
\begin{faItems}
  \item طراحی و تحلیل الگوریتم‌های حافظۀ مشترک و اولیه‌های همگام‌سازی برای سیستم‌های دارای حافظۀ پایدار (غیرفرّار).
  \item تمرکز بر اشیای بازیابی‌پذیر: ساختمان‌های دادۀ همروندی که وضعیتشان پس از خرابی سیستم سازگار می‌ماند --- مستقیماً مرتبط با سامانه‌های ذخیره‌سازی بادوام و تحمل‌پذیر خطا.
\end{faItems}

\faEntry{دستیار پژوهشی تحصیلات تکمیلی --- دانشگاه کلگری}{۲۰۲۳ تا ۲۰۲۵}
استاد راهنما: پروفسور \lr{Philipp Woelfel}
\begin{faItems}
  \item طراحی نخستین الگوریتم قفلِ بی‌انتظارِ تصادفی با فضای ثابت برای انحصار متقابل در سیستم‌های حافظۀ مشترک؛ اثبات درستی و بی‌انتظاری و تحلیل کارایی با نظریۀ احتمال.
  \item مقالۀ حاصل، \textbf{جایزۀ بهترین مقاله} را در \lr{PODC 2026}، کنفرانس شاخص محاسبات توزیع‌شده، دریافت کرد.
\end{faItems}

\faEntry{دستیار پژوهشی کارشناسی (از راه دور) --- دانشگاه جانز هاپکینز}{فوریه تا اوت ۲۰۲۳}
استاد راهنما: پروفسور \lr{Mehdi Pirooznia}
\begin{faItems}
  \item ساخت خطوط لولۀ یادگیری ماشین و جریان‌های کاری انتخاب ویژگی با \lr{Python} و \lr{R} برای پیش‌بینی ترکیب سلول‌های ایمنی در سرطان ریه؛ اجرای تحلیل بیان متفاوت و غنی‌سازی، و خودکارسازی گزارش‌گیری برای نتایج تکرارپذیر.
\end{faItems}

\end{faSection}

\begin{faSection}{پروژه‌های منتخب}

\faEntry{ساختمان‌های دادۀ همروند روی سخت‌افزار چندهسته‌ای}{۲۰۲۶}
پیاده‌سازی و ارزیابی درخت‌های جست‌وجوی دودویی همروند مبتنی بر \lr{TLE} و جدول‌های درهم‌سازی گسترش‌پذیر با حافظۀ تراکنشی محدود اینتل روی سروری \lr{NUMA} با ۱۹۲ نخ. تحلیل رفتار تثبیت و لغو تراکنش‌ها، اشتراک کاذب، اثرهای \lr{NUMA} و بازپس‌گیری حافظۀ مبتنی بر دوره با \lr{perf}.

\vspace{0.4em}
\faEntry{اپلیکیشن گفت‌وگوی امن}{\lr{\href{https://github.com/rerfan24/DNS-project}{GitHub}}}
ساخت یک اپلیکیشن گفت‌وگو با \lr{Python} در تیمی سه‌نفره، مقاوم‌شده با چند سازوکار امنیت رایانه‌ای: احراز هویت، انتقال رمزنگاری‌شده و بررسی یکپارچگی.

\end{faSection}

\begin{faSection}{تحصیلات}

\faEntry{دانشگاه واترلو}{ژانویۀ ۲۰۲۶ تاکنون}
دکتری مهندسی برق و کامپیوتر --- استاد راهنما: پروفسور \lr{Wojciech Golab}

\vspace{0.4em}
\faEntry{دانشگاه کلگری}{۲۰۲۳ تا ۲۰۲۵}
کارشناسی ارشد علوم کامپیوتر --- پایان‌نامه: \textit{«قفل بی‌انتظارِ تصادفی با فضای بهینه»}

\vspace{0.4em}
\faEntry{دانشگاه صنعتی شریف (آریامهر)}{۲۰۱۹ تا ۲۰۲۳}
کارشناسی مهندسی کامپیوتر --- پایان‌نامه: \href{https://github.com/kahbodaeini/Fraud-Detection}{\textit{«تشخیص تقلب از روی دنباله‌های تراکنش‌های بانکی»}}

\end{faSection}

\begin{faSection}{جوایز}
\faEntry{جایزۀ بهترین مقاله، چهل‌وپنجمین سمپوزیوم \lr{ACM} در مبانی محاسبات توزیع‌شده (\lr{PODC})}{۲۰۲۶}
\faEntry{جایزۀ پژوهشی دانشکده، دانشکدۀ علوم کامپیوتر، دانشگاه کلگری}{۲۰۲۳}
\faEntry{میان صد نفر برتر کشور (رتبۀ ۹۴ از حدود ۱۶۰٬۰۰۰ نفر در کنکور سراسری) و بورسیۀ بنیاد ملی نخبگان}{۲۰۱۹}
\end{faSection}

\begin{faSection}{انتشارات منتخب}
\textit{«قفل‌های بی‌انتظارِ تصادفیِ ساده و کارآمد»} --- \lr{PODC 2026}. \textbf{جایزۀ بهترین مقاله.}\\
\textit{«تحلیل خطی‌پذیری در سیستم‌های توزیع‌شدۀ نسبیتی»} --- کارگاه \lr{ApPLIED} هم‌زمان با \lr{PODC 2026}.\\
\textit{«ساخت جهانیِ سریع‌ترِ اشیای خواندن-تغییر-نوشتنِ آشکارپذیر برای حافظۀ پایدار»} --- در دست داوری، \lr{PPoPP 2027}.
\end{faSection}

\begin{faSection}{رهبری و آموزش}
\textbf{سرپرست دستیاران آموزشی}، دانشگاه کلگری (۲۰۲۳ تا ۲۰۲۶) و دانشگاه صنعتی شریف (۲۰۲۰ تا ۲۰۲۳). سرپرستی تیم‌های دستیار آموزشی تا ۲۵ نفر در درس‌های الگوریتم‌ها، ساختمان داده و مبانی نظری علوم کامپیوتر، با پشتیبانی بیش از ۴۰۰ دانشجو؛ مسئول طراحی تمرین‌ها، برگزاری کلاس‌های حل تمرین و استانداردهای ارزیابی.
\end{faSection}

\end{document}
